\begin{textsample}{POS Dim 1 – vem\_ed\_20 – Score 15.00 – vem\_ed\_20/t099.txt}  \label{ex:f1_pos_107}
Internacionalização e Intercâmbios Virtuais: tendências A mesa redonda"Novas tendências para a internacionalização do ensino \textbf{superior}", realizada no dia 15, \textbf{foi} composta pelos professores Grasiele Nascimento (Unisal/ Unesp) e José Celso Freire Junior (Unesp), com mediação de Osvaldo Succi Junior.
Grasiele \textbf{destacou} o crescimento das \textbf{ações} de Internacionalização em Casa no \textbf{contexto} da pandemia, além de ressaltar a maior acessibilidade da mobilidade virtual aos alunos, por \textbf{questões} logísticas e principalmente financeiras."Apenas 1,4 \% dos estudantes de graduação brasileiros conseguem fazer \textbf{intercâmbios} no Exterior", comentou.
José Celso Freire Junior, por outro lado, relatou \textbf{que} os Intercâmbios Virtuais na Unesp não \textbf{foram} resposta à pandemia, vieram antes."O \textbf{programa} Ciência Sem Fronteiras, até 2013, abriu espaço para a internacionalização e o \textbf{desenvolvimento} de parcerias posteriores", relatou.
A Unesp tem como base um Plano Estratégico de Internacionalização e uma política de \textbf{idiomas} \textbf{que} norteia \textbf{ações} como cursos de línguas, exames de certificação e preparação de professores.
Nos dias 14 e 15 de setembro, Ana Carolina Freschi, Neusa Haruka, Patrícia Patrício e Regiane Moreira, da equipe dos PCIs / Cesu, conduziram dois workshops.
O primeiro abordou design de projetos e o segundo tratou das \textbf{formas} de avaliação de um PCI.
Nove trabalhos apresentados na área"COIL e Intercâmbios Virtuais"\textbf{compartilharam} experiências e \textbf{reflexões} de professores envolvidos com Projetos Colaborativos Internacionais (PCIs / Cesu) nas Fatecs. \textbf{São} eles: “Abordagem COIL: a \textbf{interculturalidade} na formação do ensino \textbf{superior}”, de Vanessa Cristhina Gatto e Regiane Souza Camargo Moreira; “Ensino de línguas por meio de CLIL em uma experiência COIL”, de Ana Lúcia Leme Prestes e Edilene Gasparini Fernandes; “Intercâmbio cultural virtual: \textbf{impacto} na aprendizagem da língua inglesa e \textbf{desenvolvimento} de \textbf{habilidades} \textbf{linguísticas}”, de Tálita Suelen de Oliveira Guarino; “Escrituras narrativas audiovisuais: a aprendizagem \textbf{baseada} em projetos em PCI entre Fatec Ipiranga (Brasil) e Uniminuto (Colômbia)”, de Odenildo França Almeida; “International online course: English for Internationalization Purposes (EIP)”, de Maria Cláudia Nunes Delfino e Denise Abutbul; “O comércio internacional entre Brasil e Israel”, de Valdnea Auxiliadora de Paula Santos e Claudete Oliveira Kenvyn; “PCI: \textbf{reflexões} sobre a prática docente”, de Maria Regina Ferreira Miranda e Lilian de Souza; “PCI Fatec x SUNY Delhi: desmistificando crenças”, de Patrícia Januária da Silva Cunha Barbosa e Camila Maria da Costa Kami; “Supply value: estudo de caso do portfólio de embalagens Valgroup”, de Camila Martinelli Rocha Jacó e Zahira Gabriella Cruz Netro.

% matched lemmas: ação, basear, compartilhar, contexto, desenvolvimento, destacar, forma, habilidade, idioma, impacto, interculturalidade, intercâmbio, linguístico, programa, que, questão, reflexão, ser, superior
\end{textsample}
